This command line tool has the primary function to act as a bridge between data acquisition protocols and audio/control software.

{\bf Core Functionality}

The tool listens for data streams from the Lab Streaming Layer (LSL) -- commonly used in neuroscience, physiological monitoring, and sensor networks -- and converts that data into Open Sound Control (OSC) messages.
OSC is a standard protocol used for real-time communication between computers, synthesizers, and multimedia software (like Max/MSP, Pure Data, SuperCollider, or TASCAR).

{\bf Key Features}

\begin{itemize}
\item Multi-Stream Support: It can monitor multiple LSL streams
  simultaneously by creating a separate thread for each one.
\item Data Type Handling: It specifically looks for floating-point
  data (float32 or double64). If a stream contains other data types
  (like integers or strings), the tool will ignore that stream.
\item Real-Time Conversion: It pulls samples from LSL and immediately
  forwards them via OSC to a specified IP address and port.
\end{itemize}

{\bf Usage Examples}

\begin{enumerate}
\item Basic Conversion: Convert an LSL stream named "EEG" and send it
  to a local OSC server on port 9000.

\begin{verbatim}
lsl2osc -a EEG -u osc.udp://localhost:9000
\end{verbatim}

\item Multiple Streams: Convert two different streams ("HeartRate" and "Acceleration") simultaneously.

\begin{verbatim}
lsl2osc -a HeartRate -a Acceleration -u osc.udp://192.168.1.50:8000
\end{verbatim}

\item Custom OSC Path: Send data to a custom OSC path prefix.

\begin{verbatim}
lsl2osc -a MySensor -u osc.udp://localhost:9000 -p /sensors
\end{verbatim}
\end{enumerate}

{\bf How the OSC Message is Formatted}

When the tool sends data, it constructs the OSC path and payload as follows:

\begin{itemize}
\item OSC Path: [Prefix]/[StreamName]\\
  Example: If prefix is /lsl2osc and stream is EEG, the path is /lsl2osc/EEG.
\item Payload (Arguments):\\
  Timestamp: A double representing the LSL timestamp of the sample.\\
  Channel Data: A series of doubles representing the data from each channel in the stream.
\end{itemize}

{\bf Technical Details}

Dependencies: It relies on liblsl (for LSL) and liblo (for OSC).

Error Handling: If the network connection fails or the stream
disappears, it prints an error to stderr but attempts to keep running
(or sleeps while waiting for the stream to reappear).

Termination: It runs until you stop it (Ctrl+C).
